\section{Methods and Appendices}
\label{sec:methods}

\subsection*{A. Notation and formal definitions}
\label{sec:methods:notation}

We collect the central objects in one place.

\begin{definition}[Cluster and cluster set]
A \emph{cluster} is a finite multiset $\clusterset \subset \SphereD$
of unit-norm embeddings. The \emph{cluster set} of a corpus is a
partition $\clusterset_1, \ldots, \clusterset_K$ produced by $k$-means
on the $d$-dimensional embedding space, with $K$ chosen by the
elbow-of-inertia criterion on the corpus used for evaluation.
\end{definition}

\begin{definition}[Mean cluster distance $\dbar$]
For a cluster $\clusterset = \{x_1, \ldots, x_n\}$, define
\[
\dbar(\clusterset) \;=\; \frac{1}{\binom{n}{2}} \sum_{i<j}
\left(1 - \inner{x_i}{x_j}\right).
\]
This is the mean cosine distance between pairs in the cluster.
It is the natural scalar measure of cluster ``spread''.
\end{definition}

\begin{definition}[Effective dimension $\deff$]
\label{def:deff}
For a cluster $\clusterset$ with empirical covariance $\Sigma$,
let $\lambda_1 \geq \lambda_2 \geq \ldots \geq \lambda_d \geq 0$
be the eigenvalues of $\Sigma$. The \emph{effective dimension} is
the participation ratio of the spectrum
(Equation~\eqref{eq:deff_def}, restated here):
\[
\deff(\clusterset) \;=\; \frac{\bigl(\sum_i \lambda_i\bigr)^{2}}{\sum_i \lambda_i^{2}}
\;=\; \frac{1}{\sum_i p_i^{2}}, \qquad
p_i = \frac{\lambda_i}{\sum_j \lambda_j}.
\]
Equivalently, $\deff$ is the exponential of the R\'enyi-$2$
entropy of the normalised spectrum~\cite{royVetterli2007}; it
ranges in $[1, d]$. We use this single definition throughout the
paper for every reported $\deff$ value, and report
$\deff_{\mathrm{local}}$ as the median $\deff$ across a corpus's
clusters.
\end{definition}

\begin{definition}[Representative set $\reps$ and consolidation operator]
A \emph{consolidation operator} $\Psi: \clusterset \to \reps$ maps
a cluster to a representative set $\reps = \{r_1, \ldots, r_m\}
\subset \SphereD$ with $m \leq n$. The five operators we study:
\begin{itemize}
  \item \textbf{centroid}: $\reps = \{\bar{x}/\norm{\bar{x}}\}$ where
  $\bar{x} = \frac{1}{n}\sum x_i$.
  \item \textbf{importance-weighted}: $\reps = \{z/\norm{z}\}$ where
  $z = \sum w_i x_i$, $w_i \propto \|x_i\|$ (local density).
  \item \textbf{medoid}: $\reps = \{x_{i^\star}\}$ for
  $i^\star = \arg\min_i \sum_j d(x_i, x_j)$.
  \item \textbf{selective-prune}: $\reps$ retains the top-$p$
  highest-density points in the cluster, with $p$ the keep ratio.
  \item \textbf{GAC}: Geometry-Aware Consolidation router
  (\S\ref{sec:gac}), applies centroid, top-$p$ pruning, or medoid
  plus top-$r$ residuals depending on per-cluster spectral
  concentration.
\end{itemize}
\end{definition}

\begin{definition}[Identity retrieval error]
\label{def:id}
Given threshold $\theta \in (0,1)$ and cluster $\clusterset$ with
consolidation $\reps$, the identity retrieval error is
\[
\eps_{\mathrm{id}}(\theta, \clusterset, \reps) \;=\;
\Pr_{x \sim \mathrm{Unif}(\clusterset)}\!\left[
\max_{r \in \reps} \inner{x}{r} \; < \; \theta \right].
\]
Identity accuracy is $1 - \eps_{\mathrm{id}}$.
\end{definition}

\subsection*{B. Full proof of the duality theorem}
\label{sec:methods:proof}

We prove Theorem~\ref{thm:duality} in four steps, two of which are
stated as lemmas and proved in turn.

\begin{lemma}[Cap-volume lemma on $\SphereD$]
\label{lem:cap}
Let $x \sim \mathrm{Unif}(\SphereD)$ with $d \geq 2$, and let
$r \in \SphereD$. For any $\theta \in (0, 1)$,
\[
\Pr\!\left[\inner{x}{r} \geq \theta\right] \;\leq\; \tfrac{1}{2}
\left(1 - \theta^2\right)^{(d-1)/2}.
\]
\end{lemma}

\noindent\emph{Proof sketch.} By rotational symmetry, the event
reduces to $\{x_1 \geq \theta\}$ on $\SphereD$. The volume of a
spherical cap is classical~\cite{ball1997elementary,
milman1986asymptotic}; the stated form is the Ball tightening of the
standard $\exp(-d\theta^2/2)$ bound. $\square$

\begin{lemma}[Effective-dimension reduction]
Suppose the cluster $\clusterset$ lies in an effective-dimension
subspace: there is a projection $\Pi$ of rank $\deff$ such that
$\norm{(I - \Pi) x_i} \leq \eta$ for all $i$ (the cluster is
$\eta$-close to a $\deff$-dimensional subspace, with $\eta$ small
when the spectrum is concentrated). Then the cap-volume lemma
applies with $d$ replaced by $\deff$, up to an additive error
$O(\eta)$.
\end{lemma}

\noindent\emph{Proof sketch.} Project the cluster onto the effective
subspace; apply the cap-volume lemma inside the subspace; lift
using that inner products are preserved to additive error $O(\eta)$
by the projection bound. $\square$

\noindent\emph{Main proof.} Fix a cluster $\clusterset$, a
consolidation $\reps = \{r_1, \ldots, r_m\}$, and a threshold
$\theta$. For $x \in \clusterset$, the event
$\{\max_{r \in \reps} \inner{x}{r} < \theta\}$ is the intersection
$\bigcap_{j=1}^{m} \{\inner{x}{r_j} < \theta\}$. By the union bound,
\[
\Pr_{x}\!\left[\max_{r} \inner{x}{r} \geq \theta \right] \;\leq\;
\sum_{j=1}^{m} \Pr_{x}\!\left[\inner{x}{r_j} \geq \theta\right].
\]
Combining with the effective-dimension version of the cap-volume
lemma, and writing $\thetap$ for the modified threshold that
accounts for the cluster's mean-spread contribution (see
Corollary~\ref{cor:compression}),
\[
\Pr_{x}\!\left[\max_{r} \inner{x}{r} < \theta \right] \;\geq\;
1 \;-\; m \cdot c_1 \left(\tfrac{\thetap}{\dbar}\right)^{\deff},
\]
for a constant $c_1$ absorbing cap-volume and dimension-reduction
prefactors. Taking $m = O(1)$ and applying the Berry-Esseen rate
for the Gaussian approximation~\cite{esseen1942} to the one-sided
tail gives the uniform convergence used in the proof. The
resulting statement is Theorem~\ref{thm:duality}. $\square$

\subsection*{C. Full experimental protocol}
\label{sec:methods:protocol}

We unified the measurement pipeline across the $17{,}813$-cell
design as follows.

\paragraph{Step 1: Encoding.} For each (corpus, encoder) pair,
embed every passage with the encoder's canonical pooled output
and normalise to unit length. The six encoders are
BGE-large~\cite{xiao2023bge}, BGE-base~\cite{xiao2023bge},
E5-large~\cite{wang2022e5}, MiniLM~\cite{wang2020minilm},
MPNet~\cite{song2020mpnet}, and Nomic~\cite{nussbaum2024nomic};
SBERT~\cite{reimers2019sbert} and SimCSE~\cite{gao2021simcse}
are cited as prior-art anchors but not run in the present
measurement.

\paragraph{Step 2: Clustering.} Run $k$-means with $k$ set by the
elbow-of-inertia heuristic on the $20\%$ held-out set. Cluster
assignment is deterministic given a seed.

\paragraph{Step 3: Consolidation.} Apply each of the five operators
to produce a representative set for each cluster. Fix GAC's
thresholds at $\dbar_{\mathrm{safe}} = 0.75\,\thetap$,
$\dbar_{\mathrm{unsafe}} = 1.25\,\thetap$, $\tau_\rho = 0.30$
(see \S\ref{sec:gac}).

\paragraph{Step 4: Identity probing.} For each held-out passage
(the $20\%$ not used to build the index), compute its cosine
to every representative, rank, and record:
\begin{itemize}[leftmargin=1.8em, itemsep=2pt]
\item \emph{identity accuracy}: whether the top-$1$ representative
is in the passage's own cluster.
\item \emph{Recall@$k$}: whether the passage's cluster is covered by
the top-$k$ representatives for $k \in \{1, 10, 100\}$.
\item \emph{MRR@$20$}: reciprocal rank of the correct cluster within
the top $20$.
\item \emph{coverage@$0.80$}: fraction of representatives whose
cosine to some held-out query exceeds $0.80$.
\end{itemize}

\paragraph{Step 5: Paraphrase probing.} For E4, generate paraphrases
by held-out-template substitution on the query set; the paraphrase
paraphrase is calibrated so that the median embedding-cosine to the
original passage is $0.92$ (MS MARCO as the calibration corpus).

\paragraph{Step 6: Downstream QA.} For E7, concatenate the top-$5$
retrieved passages (recovered by cosine from the consolidated
representatives) as context for
Llama-$3.1$-$70$B-Instruct~\cite{llama31} served through
vLLM~\cite{kwon2023vllm}. Score EM/F$1$ with the SQuAD v$1.1$ text
normaliser~\cite{rajpurkar2016squad}.

\paragraph{Step 7: Statistics.} Each cell is rerun with $1$--$3$
seeds depending on experiment (E1: 1 seed per grid point; E2/E4/E6:
$3$ seeds; E5: $1$ seed; E7: $1$ seed; E9: $3$ seeds $\times$ $5$
epochs). We report mean and, where appropriate, s.e.\ from binomial
sampling variance.

\subsection*{D. Datasets and preparation}
\label{sec:methods:datasets}

\paragraph{MS MARCO.} Passage ranking corpus~\cite{nguyen2016msmarco};
$\sim 8.8$M passages with real-user search queries. We use the $2019$
TREC dev split.

\paragraph{Natural Questions.} Open-domain QA pairs over Wikipedia
passages~\cite{kwiatkowski2019nq}; we use the long-answer split,
$8{,}000$ train / $500$ query.

\paragraph{HotpotQA.} Multi-hop questions requiring $2$--$3$
supporting passages~\cite{yang2018hotpotqa}; $2{,}774$ train /
$500$ query; distractor setting.

\paragraph{PopQA.} Multi-paraphrase entity QA~\cite{mallen2023popqa};
$5{,}485$ train / $500$ query.

\paragraph{Wikipedia sections.} $\sim 20$K Wikipedia article
sections, treated as documents; used for E2/E5/E8/E9. Scales
cleanly to $1$M through additional page harvest.

\paragraph{arXiv titles.} $\sim 1$M arXiv paper titles (cs.\*
subset); used only for E2 to provide a high-$\deff$ corpus that
exercises the spread regime.

\paragraph{DRM-templated.} Synthetic paraphrase corpus generated
from the Deese-Roediger-McDermott false-memory
paradigm~\cite{roediger1995drm}. Each cluster is a DRM list
($\sim 15$ words) expanded with $20$ templates per word, yielding
high-cosine near-paraphrases. Used for E4/E6/E8 to produce
controlled tight-cluster regimes.

\subsection*{E. Reproduction guide}
\label{sec:methods:reproduce}

The repository contains a single entry script,
\texttt{scripts/run\_all.sh}, which reproduces every cell of the
$17{,}813$-cell design from scratch on a $4\times$H$100$ node
(\texttt{Llama-3.1-70B} cells require one node; all other cells run
on a single A$100$ or CPU). Expected wall-clock: $\approx 72$ hours
for the full sweep.

For partial replication, individual experiments are gated by
environment variables:
\begin{itemize}[leftmargin=1.8em, itemsep=2pt]
\item \texttt{E1\_ONLY=1}: $16{,}000$-cell synthetic grid,
$\approx 4$ hours on CPU.
\item \texttt{E5\_ONLY=1}: scale study through $1$M,
$\approx 8$ hours including the $1$M medoid runs.
\item \texttt{E7\_ONLY=1}: downstream QA at $70$B, requires
$4\times$H$100$ and $\approx 3$ hours wall-clock.
\end{itemize}

The repository pins all random seeds and serialises embeddings to
parquet so cluster assignments are bit-identical across replays.
